Accurate, real-time counting of small metallic fasteners on high-speed industrial conveyor belts remains a persistent bottleneck in manufacturing logistics. Manual visual inspection is cognitively demanding and error-prone, while conventional automated solutions are frequently hampered by specular glare, dense object overlap, and the prohibitive computational cost of high-accuracy detectors. These challenges collectively preclude reliable, always-on deployment at the resource-constrained factory edge.\\

To address this, this thesis presents a lightweight edge-AI framework for the real-time detection and tallying of metallic fasteners. A purpose-built dataset was constructed under a physical mock-up environment emulating target conveyor conditions, encompassing variable illumination, motion blur, and densely clustered objects. To establish a rigorous baseline and isolate architectural limitations in small-object detection, four detector configurations were benchmarked: YOLO11s, YOLO11m, YOLOv26m, and RT-DETR as a high-accuracy reference.\\

Guided by baseline limitations, this work proposes YOLO11-GED, a custom lightweight detector optimized for edge deployment. The architecture introduces the novel EMAc2f module, which synergizes Efficient Multi-Scale Attention (EMA) for occlusion-aware feature refinement with GhostBottleneck operations for parameter-efficient feature reuse. Furthermore, DySample dynamic upsampling is strategically integrated into the network neck to preserve high-frequency spatial details during feature fusion, directly mitigating the blurring of fine-grained metallic textures.\\

The optimized model is deployed on an NVIDIA Jetson Orin Nano using TensorRT acceleration, coupled with ByteTrack for temporal consistency and geometric counting logic to prevent double-counting. Experimental results demonstrate that YOLO11-GED achieves exceptional real-time performance, sustaining $89.9–100.8 FPS$ with latency under $11.5 ms$ and low hardware utilization $(GPU \approx 37–40 \%, CPU \approx 8 \%)$. Crucially, it maintains high counting accuracy in high-density, occluded scenarios where heavier baselines like RT-DETR significantly degrade.\\

Beyond the immediate industrial application, this work contributes a reproducible design-to-deployment methodology for small-object detection under strict edge-compute budgets. Validated through a successful real-world Proof of Concept (PoC), it provides empirical evidence that targeted, domain-specific architectural modifications can significantly improve small-object recall and robustness without sacrificing real-time throughput.